\begin{example}\label{ex:p142-infinite-ideal-value}
On $X=\mathbf P^1$, choose a point $p$, a smooth Hermitian metric on
$\mathcal O(p)$, and its canonical section $s_p$. Fix $\lambda>0$ and
put $\psi=\lambda\log|s_p|^2$. The family $\Gamma_\tau=\psi$ for
$\tau<c$ satisfies the elementary hypotheses of
\cref{thm:p142-raw-transform-usc}; no assertion that this particular
representative is an $\mathcal I$-model potential is needed.

Writing $\mathfrak m_p$ for the ideal sheaf of $p$, one has
\[
       \mathcal I(m\psi)=\mathfrak m_p^{\lfloor m\lambda\rfloor}.
\]
Indeed, near $p$ a coordinate $z$ expresses $\psi$ as
$\lambda\log|z|^2$ plus a smooth function. A germ of order $q$ is
integrable with the weight $\mathrm e^{-m\psi}$ exactly when
$q-m\lambda>-1$, which is equivalent to $q\geq\lfloor m\lambda\rfloor$.
Away from $p$ the ideal is the unit ideal. It follows that
\[
 \nu_v^{\mathcal I}(\psi)=\lambda v(\mathfrak m_p),
 \qquad F_\Gamma(v)=c-\lambda v(\mathfrak m_p).
\]
At a semivaluation with $v(\mathfrak m_p)=+\infty$, each sufficiently
large multiplier-ideal value is infinite, so the displayed formulas give
$F_\Gamma(v)=-\infty$. For the smaller integers with
$\lfloor m\lambda\rfloor=0$, the ideal is the unit ideal and its value
is zero; no expression $0\cdot(+\infty)$ is used. Thus the infinite-value
case in the theorem records actual support information and cannot be
discarded by restricting to divisorial valuations.
\end{example}
